\section{Related Work}
\label{sec:related}

The project sits at the intersection of three lines of research:
\emph{generative modelling} through diffusion and flow matching, the
\emph{diffusion transformer} architecture that has become the workhorse of
modern generative models, and \emph{world models} for autonomous driving and
navigation. Figure~\ref{fig:taxonomy} organizes the most relevant works along
these axes. We then treat each line in turn: Section~\ref{sec:diffusion-theory}
develops the theory of diffusion and flow matching, Section~\ref{sec:dit}
describes the diffusion transformer, and Section~\ref{sec:world-models} reviews
world models from the general concept through to the driving and navigation
models that motivate our probe.

\begin{figure}[t]
\centering
\begin{tikzpicture}[
  font=\small,
  grow=down,
  level distance=12mm,
  root/.style={draw, rounded corners, fill=gray!15, align=center, inner sep=4pt},
  cat/.style={draw, rounded corners, fill=blue!10, align=center, inner sep=3pt},
  leaf/.style={draw, rounded corners, fill=green!10, align=center, inner sep=3pt, font=\scriptsize},
  ar/.style={-{Stealth[length=2mm]}, semithick},
  >={Stealth[length=2mm]}
]
\node[root] (root) at (0,0) {Generative \& predictive models\\for driving};

% three category nodes
\node[cat] (gen) at (-5.3,-1.4) {Generative\\objectives};
\node[cat] (arch) at (0,-1.4) {Architecture};
\node[cat] (wm) at (5.3,-1.4) {World models};

\draw[ar] (root) -- (gen);
\draw[ar] (root) -- (arch);
\draw[ar] (root) -- (wm);

% generative leaves
\node[leaf] (diff) at (-6.6,-3.0) {Diffusion\\(DDPM)~\cite{ho_ddpm}};
\node[leaf] (fm) at (-4.5,-3.0) {Flow matching /\\rectified flow~\cite{lipman_fm}};
\node[leaf] (disc) at (-6.6,-4.4) {Masked / discrete\\token models~\cite{maskgit}};
\node[leaf] (vae) at (-4.5,-4.4) {(VQ-)VAE\\tokenizers~\cite{vqgan}};
\draw[ar] (gen) -- (diff);
\draw[ar] (gen) -- (fm);
\draw[ar] (gen) -- (disc);
\draw[ar] (gen) -- (vae);

% architecture leaves
\node[leaf] (unet) at (0,-3.0) {U-Net\\backbones};
\node[leaf] (dit) at (0,-4.1) {Diffusion\\transformer (DiT)~\cite{dit}};
\node[leaf] (stdit) at (0,-5.2) {Spatio-temporal\\DiT / CDiT~\cite{nwm}};
\draw[ar] (arch) -- (unet);
\draw[ar] (arch) -- (dit);
\draw[ar] (arch) -- (stdit);

% world model leaves
\node[leaf] (drive) at (6.7,-3.0) {Driving WMs:\\GAIA, Vista,\\GEM};
\node[leaf] (orbis) at (4.3,-3.0) {\textbf{Orbis}~\cite{orbis}};
\node[leaf] (nav) at (6.7,-4.6) {Navigation\\WM (NWM)~\cite{nwm}};
\node[leaf] (game) at (4.3,-4.6) {Game / sim WMs:\\DIAMOND,\\GameNGen};
\draw[ar] (wm) -- (drive);
\draw[ar] (wm) -- (orbis);
\draw[ar] (wm) -- (nav);
\draw[ar] (wm) -- (game);

% cross link: Orbis combines flow matching (a generative objective) with the
% spatio-temporal DiT (an architecture); routed below the tree to avoid nodes
\draw[ar, dashed, red!60] (orbis.south) .. controls (3.1,-3.6) and (2.3,-4.6) .. (stdit.east);
\draw[ar, dashed, red!60] (orbis.south) .. controls (2.9,-4.2) and (-3.6,-7.2) .. (fm.south);
\end{tikzpicture}
\caption{A taxonomy of the works reviewed in this section. Driving world models
can be classified by their \emph{generative objective} (continuous diffusion /
flow matching vs.\ discrete masked-token modelling), their \emph{architecture}
(convolutional U-Net vs.\ transformer-based DiT variants), and their
\emph{role} (open-world driving generation, navigation planning, or game/sim
emulation). Dashed red edges highlight that Orbis combines a flow-matching
objective with a spatio-temporal diffusion transformer.}
\label{fig:taxonomy}
\end{figure}

\subsection{Diffusion and flow matching}
\label{sec:diffusion-theory}

\paragraph{Denoising diffusion probabilistic models.}
Diffusion models~\cite{ho_ddpm} define a generative process by reversing a
fixed forward process that gradually corrupts data into noise. Given a data
sample $\mathbf{x}_0$, the forward process produces a sequence of increasingly
noisy variables $\mathbf{x}_1,\dots,\mathbf{x}_T$ according to a Gaussian
transition with a predefined variance schedule $\{\alpha_t\}$. A convenient
property is that any noise level can be sampled in closed form,
\begin{equation}
\mathbf{x}_t = \sqrt{\bar{\alpha}_t}\,\mathbf{x}_0
+ \sqrt{1-\bar{\alpha}_t}\,\boldsymbol{\epsilon},
\qquad \boldsymbol{\epsilon}\sim\mathcal{N}(0,I),
\end{equation}
where $\bar{\alpha}_t=\prod_{s\le t}\alpha_s$. The model learns the reverse
process $p_\theta(\mathbf{x}_{t-1}\mid\mathbf{x}_t)$, which gradually denoises a
pure-noise sample back to data. Training reduces to a simple regression in
which a network $\boldsymbol{\epsilon}_\theta$ predicts the noise that was added,
\begin{equation}
\mathcal{L}_{\text{simple}}(\theta)
= \mathbb{E}_{t,\mathbf{x}_0,\boldsymbol{\epsilon}}
\big[\,\lVert \boldsymbol{\epsilon}_\theta(\mathbf{x}_t,t)
- \boldsymbol{\epsilon} \rVert_2^2\,\big].
\end{equation}
New samples are generated by initializing
$\mathbf{x}_T\sim\mathcal{N}(0,I)$ and iterating the learned reverse
transitions. \emph{Classifier-free guidance}~\cite{ho_ddpm} is the standard
mechanism for conditional generation: the model is trained both with and
without the conditioning $c$ (the latter by dropping $c$ for a learned null
embedding), and at sampling time the two predictions are extrapolated,
$\hat{\boldsymbol{\epsilon}}(\mathbf{x}_t,c)
=\boldsymbol{\epsilon}_\theta(\mathbf{x}_t,\varnothing)
+ s\big(\boldsymbol{\epsilon}_\theta(\mathbf{x}_t,c)
-\boldsymbol{\epsilon}_\theta(\mathbf{x}_t,\varnothing)\big)$,
with guidance scale $s>1$ trading diversity for fidelity to the condition.

\paragraph{Flow matching.}
Flow matching~\cite{lipman_fm} reframes generation as learning a continuous
\emph{velocity field} that transports a simple prior distribution to the data
distribution along a probability path, rather than reversing a discrete
noising chain. The simplest and most widely used instance defines a linear
(``rectified'') path between a data point $\mathbf{x}$ and a noise sample
$\boldsymbol{\epsilon}$,
\begin{equation}
\mathbf{x}^{\tau} = (1-\tau)\,\mathbf{x} + \tau\,\boldsymbol{\epsilon},
\qquad \tau\in[0,1],
\end{equation}
whose time-derivative is the constant target velocity
$\tfrac{d\mathbf{x}^{\tau}}{d\tau}=\boldsymbol{\epsilon}-\mathbf{x}$. The model
$\mathbf{v}_\theta$ is trained to regress this velocity,
\begin{equation}
\mathcal{L}_{\text{FM}}(\theta)
= \mathbb{E}_{\tau,\mathbf{x},\boldsymbol{\epsilon}}
\big[\,\lVert \mathbf{v}_\theta(\mathbf{x}^{\tau},\tau)
- (\boldsymbol{\epsilon}-\mathbf{x}) \rVert_2^2\,\big],
\label{eq:fm-loss}
\end{equation}
and samples are produced by integrating the ordinary differential equation
$\tfrac{d\mathbf{x}^{\tau}}{d\tau}=\mathbf{v}_\theta(\mathbf{x}^{\tau},\tau)$
from the prior at $\tau=1$ back to the data manifold at $\tau=0$. Flow matching
can be seen as a continuous-time generalization of diffusion: it shares the
same denoising intuition but offers straighter transport paths, which permit
fewer integration steps and tend to be more stable to train. This is the
objective Orbis uses for its continuous world model
(Equations~\eqref{eq:fm-forward-intro}--\eqref{eq:fm-loss-intro}), and it is a
natural choice for the trajectory-generation head of our probe because the 2D
trajectory space is low-dimensional and benefits from the short sampling
schedules that flow matching enables.

\paragraph{Continuous vs.\ discrete generation.}
A recurring design axis, made explicit by Orbis, is whether to model a
continuous latent (with diffusion or flow matching) or a discrete sequence of
codebook tokens (with autoregressive or masked generative
modelling~\cite{maskgit}). Discrete token models inherit the machinery of large
language models and can be attractive for long rollouts, but Orbis finds that
on driving video the continuous, flow-matching formulation is less brittle to
design choices and yields markedly better long-horizon predictions. Both
families rely on a (VQ-)VAE-style tokenizer~\cite{vqgan} to map pixels to a
compact latent; the hybrid Orbis tokenizer supports both at once.

\subsection{The diffusion transformer}
\label{sec:dit}

Early diffusion models used convolutional U-Net backbones. The Diffusion
Transformer (DiT)~\cite{dit} showed that the U-Net inductive bias is not
essential: a standard transformer operating on latent patches matches or
exceeds U-Nets and, importantly, \emph{scales} predictably --- sample quality
(FID) improves smoothly as the model's compute (Gflops), through either depth,
width, or token count, increases. DiT operates in the latent space of a
pretrained VAE: the spatial latent is ``patchified'' into a sequence of tokens,
processed by a stack of transformer blocks, and decoded back into a noise (and
covariance) prediction. The key design finding is how to inject conditioning
--- the diffusion timestep $t$ and any class or context signal $c$. Among
in-context tokens, cross-attention, and adaptive layer normalization (AdaLN),
the \emph{adaLN-Zero} variant works best: the conditioning vector regresses
per-block scale and shift parameters $(\gamma,\beta)$ for the layer norms and an
additional residual scale $\alpha$ that is zero-initialized so that each block
starts as the identity. This is precisely the conditioning mechanism Orbis
adopts to inject the diffusion time and the optional ego-motion signal, and it
is the most compute-efficient of the three.

Figure~\ref{fig:dit} sketches the block at a general level. A token sequence
enters; AdaLN modulates a multi-head self-attention sub-layer and then a
feed-forward (MLP) sub-layer; residual connections, scaled by $\alpha$, wrap
each. Stacking $N$ such blocks and conditioning every block on $(t,c)$ gives the
full backbone. Driving and navigation world models extend this design along the
\emph{temporal} axis. Orbis interleaves spatial and temporal attention within
each block (an STDiT), so that a token attends both to other patches in the
same frame and to the same region across frames. The navigation world model
NWM~\cite{nwm} instead proposes a Conditional Diffusion Transformer (CDiT) that
restricts expensive self-attention to the target frame and uses cross-attention
to read context frames, making the cost \emph{linear} in the number of context
frames rather than quadratic --- a choice that lets it scale to a billion
parameters. The general-level scheme in Figure~\ref{fig:dit} is the common
ancestor of all of these variants.

\begin{figure}[t]
\centering
\begin{tikzpicture}[
  font=\small,
  box/.style={draw, rounded corners, align=center, inner sep=3pt, minimum height=7mm},
  op/.style={box, fill=blue!10},
  attn/.style={box, fill=orange!18},
  mlp/.style={box, fill=green!14},
  norm/.style={box, fill=gray!12},
  cond/.style={box, fill=red!10},
  add/.style={draw, circle, inner sep=1pt, fill=white},
  ar/.style={-{Stealth[length=2mm]}, semithick},
  >={Stealth[length=2mm]}
]
% main column
\node[box, fill=gray!8] (in) at (0,0) {Input tokens\\(patchified latent)};
\node[norm] (ln1) at (0,1.3) {Layer Norm};
\node[attn] (mha) at (0,2.7) {Multi-Head\\Self-Attention};
\node[add] (a1) at (0,4.0) {$+$};
\node[norm] (ln2) at (0,5.3) {Layer Norm};
\node[mlp] (mlpb) at (0,6.7) {Feed-Forward\\(MLP)};
\node[add] (a2) at (0,8.0) {$+$};
\node[box, fill=gray!8] (out) at (0,9.3) {Output tokens};

\draw[ar] (in) -- (ln1);
\draw[ar] (ln1) -- (mha);
\draw[ar] (mha) -- (a1);
\draw[ar] (a1) -- (ln2);
\draw[ar] (ln2) -- (mlpb);
\draw[ar] (mlpb) -- (a2);
\draw[ar] (a2) -- (out);
% residuals
\draw[ar] (in.east) -- ++(1.7,0) |- (a1.east) node[pos=0.07, above, font=\scriptsize] {residual};
\draw[ar] (a1.west) -- ++(-1.7,0) |- (a2.west) node[pos=0.25, left, font=\scriptsize] {residual};

% conditioning column
\node[cond, align=center] (cond) at (4.6,2.7) {Conditioning\\$t$ (diffusion step)\\$c$ (context / action)};
\node[cond] (mlp2) at (4.6,4.6) {MLP};
\draw[ar] (cond) -- (mlp2);
% adaln scale/shift/scale arrows: fan out from the MLP to the left column
\draw[ar, dashed, red!70] (mlp2.west) .. controls (3.0,4.7) and (2.0,5.3) .. (ln2.east)
  node[pos=0.62, above, font=\scriptsize] {$\gamma_2,\beta_2$};
\draw[ar, dashed, red!70] (mlp2.west) -- (a1.east)
  node[pos=0.55, above, font=\scriptsize] {$\alpha_1$};
\draw[ar, dashed, red!70] (mlp2.west) .. controls (2.7,3.1) and (1.9,1.6) .. (ln1.east)
  node[pos=0.62, right, font=\scriptsize] {$\gamma_1,\beta_1$};
\node[font=\scriptsize, align=center, red!70] at (4.6,1.3) {AdaLN-Zero:\\regress $(\gamma,\beta,\alpha)$};

\node[font=\small\bfseries, align=center] at (0,10.2) {DiT block ($\times N$)};
\end{tikzpicture}
\caption{The diffusion transformer (DiT) block at a general level. Patchified
latent tokens pass through a multi-head self-attention sub-layer and a
feed-forward (MLP) sub-layer, each wrapped in a residual connection. The
diffusion timestep $t$ and the conditioning $c$ are injected through adaptive
layer normalization: a small MLP regresses per-sub-layer scale/shift
$(\gamma,\beta)$ and a residual scale $\alpha$ (zero-initialized, ``adaLN-Zero'').
Stacking $N$ blocks gives the backbone; driving models such as Orbis add
temporal attention (STDiT) and navigation models such as NWM use a
context cross-attention (CDiT) variant. This scheme is intentionally
schematic --- the detailed per-block diagram is left to be refined separately.}
\label{fig:dit}
\end{figure}

\subsection{World models}
\label{sec:world-models}

\paragraph{The world-model concept.}
A world model is a learned model of an environment's dynamics: given the
current state (or a history of observations) and optionally an action, it
predicts the next state, and possibly a reward. The motivation is that an agent
equipped with such a model can ``imagine'' the outcome of candidate actions
internally, which supports sample-efficient reinforcement learning,
representation learning, and planning. Historically, world models were studied
in games and other simulated environments, where the dynamics are clean and
unlimited interaction is available. The recent surge of high-quality video
generative models has made it feasible to learn world models directly from raw
video of the real world, turning the predictive model into a learned simulator.
In the driving domain this is especially attractive because raw front-camera
video is abundant and unlabeled, so a video world model can sidestep much of the
manual perception engineering that conventional stacks require.

\paragraph{Popular video and driving world models.}
Several lines of work have established video prediction as a viable substrate
for world modelling. In games and simulation, DIAMOND and GameNGen use
diffusion models to emulate game engines such as Atari and Doom from frames and
actions. For driving, GAIA-1 was an influential early demonstration of
controllable future generation conditioned on action and text, using a discrete
token model. Diffusion-based driving models including Vista and GEM fine-tune
general-purpose pretrained video diffusion backbones (such as Stable Video
Diffusion) to produce high-resolution, high-frame-rate rollouts of up to around
fifteen seconds, and DrivingWorld further extends rollout coherence. Orbis
positions itself critically against this trend: it observes that these large,
pretrained-backbone models degrade quickly on long horizons and on demanding
maneuvers such as turns, often stopping prematurely or drifting into unrealistic
paths. To quantify this, Orbis introduces a distribution-level, trajectory-based
evaluation --- mapping generated videos to pose sequences and measuring realism
and coverage with precision/recall, Fréchet distance, and average displacement
error --- since standard video metrics like FVD are insensitive to scene
dynamics and ego-motion. This evaluation philosophy is directly relevant to our
probe, which likewise assesses the \emph{distribution} of trajectories rather
than a single point estimate.

\paragraph{Orbis.}
Orbis~\cite{orbis} (described in detail in Section~\ref{sec:orbis}) contributes
along two axes that matter here. First, it is trained from raw video only and is
far more economical than its competitors --- $469$M parameters and $280$ hours
of video versus the thousands of hours and larger backbones used elsewhere ---
which suggests the learned latent is an efficient, general representation rather
than a memorized one. Second, its hybrid tokenizer enables a controlled
comparison between continuous and discrete world models, concluding in favour of
the continuous flow-matching model. For this report, the relevant consequence is
that Orbis offers a compact, well-trained, raw-video latent whose informativeness
for downstream driving tasks is an open and interesting question --- precisely
what the attentive diffusion probe is designed to measure.

\paragraph{Navigation world models.}
The Navigation World Model (NWM)~\cite{nwm} makes the planning use of a video
world model explicit, which is conceptually the closest precedent to our study.
NWM is a controllable video generation model trained on egocentric video of
human and robotic agents together with their navigation actions
$(\Delta x,\Delta y,\Delta\phi)$ and a time shift. It uses the CDiT architecture
mentioned in Section~\ref{sec:dit} to autoregressively predict future latent
states given past states and actions. Crucially, after training, NWM is used for
\emph{planning}: it simulates candidate action sequences and scores how well the
resulting imagined future reaches a goal, either optimizing actions from scratch
in a model-predictive-control loop or ranking trajectories proposed by an
external policy. This demonstrates the broader thesis behind our work --- that
the representation and rollouts of a video world model carry actionable
information about future motion. Where NWM \emph{generates and scores} full
future videos to plan, our probe takes the complementary, lighter-weight route:
it reads the world model's latent directly and learns a small diffusion head to
predict the future 2D trajectory, testing how much of that planning-relevant
information is already linearly accessible in the frozen Orbis latent.
